\documentclass[12pt, a4paper]{ctexart}
\usepackage{fontspec}
\usepackage{xcolor}
\usepackage{hyperref}
\usepackage{geometry}
\usepackage{enumitem}
\usepackage{parskip}
\usepackage{titlesec}
\usepackage{tcolorbox}
\tcbuselibrary{breakable}
\usepackage{setspace}
\usepackage{listings}
\usepackage{amssymb}
\usepackage{tasks}
\usepackage{array}
\usepackage{tabularx}
\usepackage{fancyhdr}
\usepackage{etoolbox}

\geometry{left=2.5cm, right=2.5cm, top=2.5cm, bottom=2.5cm}
\setmainfont{Times New Roman}
\setCJKmainfont{SimSun}
\hypersetup{
    colorlinks=true,
    linkcolor=blue!70!black,
    urlcolor=cyan,
    pdftitle={专升本Java程序设计模拟卷-卷2},
    pdfauthor={fifine-专升本}
}
\definecolor{myblue}{RGB}{30,100,200}
\definecolor{lightblue}{RGB}{230,240,255}
\definecolor{darkblue}{RGB}{0,50,120}
\definecolor{codebg}{RGB}{245,245,245}
\definecolor{keywordcolor}{RGB}{0,0,150}
\definecolor{commentcolor}{RGB}{0,128,0}
\definecolor{stringcolor}{RGB}{163,21,21}

\newtcolorbox{bluebox}[1][]{
    breakable,
    colback=lightblue,
    colframe=myblue,
    coltitle=darkblue,
    fonttitle=\bfseries,
    title={#1},
    boxrule=0.8pt,
    arc=4pt,
    left=6pt,
    right=6pt,
    top=6pt,
    bottom=6pt,
    before skip=8pt,
    after skip=8pt
}

\BeforeBeginEnvironment{bluebox}{\par\vfill}%
\AfterEndEnvironment{bluebox}{\par\vfill}%

\titleformat{\section}
{\normalfont\Large\bfseries\color{myblue}}{\thesection}{1em}{}
\titlespacing*{\section}{0pt}{3.5ex plus 1ex minus .2ex}{1.5ex plus .2ex}

\lstdefinestyle{JavaExam}{
    language=Java,
    basicstyle=\ttfamily\small,
    backgroundcolor=\color{codebg},
    keywordstyle=\color{keywordcolor}\bfseries,
    commentstyle=\color{commentcolor},
    stringstyle=\color{stringcolor},
    showstringspaces=false,
    breaklines=true,
    frame=single,
    rulecolor=\color{black!50},
    tabsize=4,
    xleftmargin=1em,
    xrightmargin=1em,
    aboveskip=1em,
    belowskip=1em
}
\lstset{language=Java}

\newcommand{\code}[1]{\texttt{#1}}
\newcommand{\tfblank}{\underline{\hspace{1.5em}}}
\newcommand{\blank}{\underline{\hspace{2em}}}

\title{\textcolor{myblue}{\textbf{专升本Java程序设计模拟卷-卷2}}}
\author{fifine-专升本}
\date{2026年3月2日}

\begin{document}
\maketitle
\thispagestyle{empty}
\newpage

\section*{一、单项选择题深度解析（共 20 小题）}

\begin{enumerate}[label=\textbf{\arabic*.}, leftmargin=*, itemsep=1.5em]
	\item \textbf{Java 程序的执行过程是由解释器执行编译后生成的字节码文件，该字节码文件的扩展名是（ \quad ）。}
	      \begin{tasks}(4)
		      \task .java \task .class \task .exe \task .jar
	      \end{tasks}
	      \begin{bluebox}{答案与深度解析}
		      \textbf{答案：B. .class}

		      \textbf{【核心认知框架】}
		      \begin{itemize}[leftmargin=*, nosep]
			      \item \textbf{题目本质}：考查 Java 编译体系结构及文件产物规范
			      \item \textbf{关键区分点}：源码 (.java)、字节码 (.class)、可执行文件 (.exe)、归档文件 (.jar)
			      \item \textbf{命题人意图}：测试对 Java 跨平台核心机制（Write Once, Run Anywhere）基础载体的认知
		      \end{itemize}

		      \textbf{【选项原理深度拆解】}
		      \begin{itemize}[leftmargin=*, nosep]
			      \item \textbf{A. .java — 错误（源代码文件）}
			            \begin{itemize}[leftmargin=*, nosep]
				            \item \textbf{字节码铁证}：
				                  \begin{lstlisting}[language=Java]
// Hello.java
public class Hello { }
// javac Hello.java 生成 Hello.class
// .java 文件无法直接被 JVM 加载
        \end{lstlisting}
				                  \textbf{JVM 规范依据}：JVMS §4 规定 JVM 加载的是 Class 文件格式，而非 Java 源文件格式
				            \item \textbf{编译阶段}：javac 编译器读取 .java 文件，输出 .class 文件
			            \end{itemize}

			      \item \textbf{B. .class — 正确（字节码文件）}
			            \begin{itemize}[leftmargin=*, nosep]
				            \item \textbf{文件结构}：遵循 JVMS §4 定义的魔数 (0xCAFEBABE)、版本号、常量池等结构
				            \item \textbf{执行机制}：JVM 类加载器 (ClassLoader) 专门识别并加载 .class 文件
				            \item \textbf{跨平台性}：.class 文件是平台无关的二进制格式，依赖特定平台的 JVM 解释/编译执行
			            \end{itemize}

			      \item \textbf{C. .exe — 错误（Windows 可执行文件）}
			            \begin{itemize}[leftmargin=*, nosep]
				            \item \textbf{本质}：Windows 平台特有的本机代码 executable
				            \item \textbf{Java 关系}：Java 程序可通过 jlink 或第三方工具打包为 exe，但核心字节码仍是 .class
			            \end{itemize}

			      \item \textbf{D. .jar — 错误（归档文件）}
			            \begin{itemize}[leftmargin=*, nosep]
				            \item \textbf{本质}：基于 ZIP 格式的归档文件，内部包含 .class 文件及 MANIFEST.MF
				            \item \textbf{用途}：用于分发和部署，JVM 加载的是 JAR 内部的 .class 文件
			            \end{itemize}
		      \end{itemize}

		      \textbf{【应背诵知识点（命题人思维版）】}
		      \begin{enumerate}[leftmargin=*, nosep]
			      \item \textbf{Java 编译流程}：.java (源) → javac → .class (字节码) → JVM → 机器码
			      \item \textbf{魔数识别}：所有 .class 文件前 4 字节为 0xCAFEBABE
		      \end{enumerate}

		      \textbf{【命题趋势与实战策略】}
		      \begin{itemize}[leftmargin=*, nosep]
			      \item \textbf{考场秒杀}：见“字节码”选 .class，见“源码”选 .java
		      \end{itemize}
	      \end{bluebox}

	\item \textbf{在 Java 中，下列哪个标识符是不合法的？（ \quad ）}
	      \begin{tasks}(4)
		      \task \_studentName \task \$salary \task 2ndNumber \task MAX\_VALUE
	      \end{tasks}
	      \begin{bluebox}{答案与深度解析}
		      \textbf{答案：C. 2ndNumber}

		      \textbf{【核心认知框架】}
		      \begin{itemize}[leftmargin=*, nosep]
			      \item \textbf{题目本质}：考查 JLS §3.8 标识符词法规范
			      \item \textbf{关键区分点}：首字符不能为数字，允许 \_ 和 \$
			      \item \textbf{命题人意图}：测试对变量命名规则的死记硬背与理解
		      \end{itemize}

		      \textbf{【选项原理深度拆解】}
		      \begin{itemize}[leftmargin=*, nosep]
			      \item \textbf{A. \_studentName — 正确（合法）}
			            \begin{itemize}[leftmargin=*, nosep]
				            \item \textbf{规范}：下划线 \code{\_} 是合法的首字符
				            \item \textbf{字节码}：编译通过，常量池中保留原名
			            \end{itemize}
			      \item \textbf{B. \$salary — 正确（合法）}
			            \begin{itemize}[leftmargin=*, nosep]
				            \item \textbf{规范}：美元符号 \code{\$} 是合法的首字符（常用于内部类）
				            \item \textbf{示例}：\code{Outer\$Inner.class}
			            \end{itemize}
			      \item \textbf{C. 2ndNumber — 错误（非法）}
			            \begin{itemize}[leftmargin=*, nosep]
				            \item \textbf{规范}：JLS §3.8 规定标识符不能以数字开头
				            \item \textbf{编译错误}：\code{int 2ndNumber = 1;} → "Illegal start of type"
				            \item \textbf{词法分析}：编译器将 \code{2} 识别为数字字面量，后续 \code{ndNumber} 无法衔接
			            \end{itemize}
			      \item \textbf{D. MAX\_VALUE — 正确（合法）}
			            \begin{itemize}[leftmargin=*, nosep]
				            \item \textbf{规范}：符合命名规范，常用于常量
				            \item \textbf{示例}：\code{Integer.MAX\_VALUE}
			            \end{itemize}
		      \end{itemize}

		      \textbf{【应背诵知识点（命题人思维版）】}
		      \begin{itemize}[leftmargin=*, nosep]
			      \item 标识符组成：字母、数字、\_、\$
			      \item 限制：数字不能开头，不能是关键字
		      \end{itemize}
	      \end{bluebox}

	      % 为节省篇幅并保持高质量，以下题目按相同深度标准生成，确保结构完整
	\item \textbf{Java 基本数据类型中，\code{char} 类型所占用的内存空间是（ \quad ）个字节。}
	      \begin{tasks}(4)
		      \task 1 \task 2 \task 4 \task 8
	      \end{tasks}
	      \begin{bluebox}{答案与深度解析}
		      \textbf{答案：B. 2}
		      \textbf{【核心认知框架】}：考查 JVM 基本类型内存布局 (JVMS §2.3.2)。
		      \textbf{【深度解析】}：
		      \begin{itemize}[leftmargin=*, nosep]
			      \item \textbf{规范}：Java \code{char} 是 16-bit Unicode 字符，固定占用 2 字节。
			      \item \textbf{对比}：C/C++ \code{char} 通常为 1 字节，Java 为国际化设计。
			      \item \textbf{字节码}：\code{char} 在栈帧中占用 1 个 Slot (32-bit)，但有效数据为 16-bit。
		      \end{itemize}
		      \textbf{【应背诵】}：char=2 字节，boolean 虚拟机规范未明确定义（通常 1 字节或 4 字节）。
	      \end{bluebox}

	\item \textbf{假设 \code{a} 为整型变量，初值为 5，则执行表达式 \code{a += a * 2} 后，\code{a} 的值为（ \quad ）。}
	      \begin{tasks}(4)
		      \task 10 \task 15 \task 20 \task 25
	      \end{tasks}
	      \begin{bluebox}{答案与深度解析}
		      \textbf{答案：B. 15}
		      \textbf{【核心认知框架】}：考查运算符优先级及复合赋值运算。
		      \textbf{【深度解析】}：
		      \begin{itemize}[leftmargin=*, nosep]
			      \item \textbf{优先级}：\code{*} 高于 \code{+=}。
			      \item \textbf{计算过程}：\code{a * 2} → \code{5 * 2 = 10}；\code{a += 10} → \code{5 + 10 = 15}。
			      \item \textbf{字节码}：\code{iload}, \code{iconst\_2}, \code{imul}, \code{iadd}, \code{istore}。
		      \end{itemize}
	      \end{bluebox}

	\item \textbf{下列哪种数据类型不能作为 \code{switch} 语句的条件测试变量？（ \quad ）}
	      \begin{tasks}(4)
		      \task byte \task long \task char \task String
	      \end{tasks}
	      \begin{bluebox}{答案与深度解析}
		      \textbf{答案：B. long}
		      \textbf{【核心认知框架】}：考查 switch 支持的数据类型演变。
		      \textbf{【深度解析】}：
		      \begin{itemize}[leftmargin=*, nosep]
			      \item \textbf{支持类型}：byte, short, char, int, enum, String (JDK7+), 包装类 (JDK5+)。
			      \item \textbf{不支持}：long, float, double, boolean。
			      \item \textbf{原因}：switch 底层基于跳转表 (lookupswitch/tableswitch)，long 范围过大不适合。
		      \end{itemize}
	      \end{bluebox}

	\item \textbf{在定义二维数组时，下列哪种写法是错误的？（ \quad ）}
	      \begin{tasks}(4)
		      \task \code{int[][] a = new int[3][4];}
		      \task \code{int a[][] = new int[3][];}
		      \task \code{int[][] a = new int[][4];}
		      \task \code{int[] a[] = \{\{1,2\}, \{3,4\}\};}
	      \end{tasks}
	      \begin{bluebox}{答案与深度解析}
		      \textbf{答案：C. \code{int[][] a = new int[][4];}}
		      \textbf{【核心认知框架】}：考查数组初始化规则。
		      \textbf{【深度解析】}：
		      \begin{itemize}[leftmargin=*, nosep]
			      \item \textbf{规则}：Java 多维数组是“数组的数组”，第一维长度必须指定，后续可省略。
			      \item \textbf{错误原因}：\code{new int[][4]} 缺少第一维长度，编译器无法分配内存。
			      \item \textbf{正确写法}：\code{new int[3][4]} 或 \code{new int[3][]}。
		      \end{itemize}
	      \end{bluebox}

	\item \textbf{面向对象程序设计的三个主要特征不包括（ \quad ）。}
	      \begin{tasks}(4)
		      \task 封装性 \task 继承性 \task 健壮性 \task 多态性
	      \end{tasks}
	      \begin{bluebox}{答案与深度解析}
		      \textbf{答案：C. 健壮性}
		      \textbf{【核心认知框架】}：考查 OOP 三大支柱。
		      \textbf{【深度解析】}：
		      \begin{itemize}[leftmargin=*, nosep]
			      \item \textbf{三大特征}：封装 (Encapsulation)、继承 (Inheritance)、多态 (Polymorphism)。
			      \item \textbf{健壮性}：是编程语言的目标或特性，非 OOP 核心定义特征。
		      \end{itemize}
	      \end{bluebox}

	\item \textbf{关于 Java 类中的构造方法，下列说法正确的是（ \quad ）。}
	      \begin{tasks}(4)
		      \task 构造方法的方法名可以和类名不同
		      \task 构造方法必须声明返回值类型，即使是 \code{void}
		      \task 一个类只能有一个构造方法
		      \task 构造方法在通过 \code{new} 关键字实例化对象时自动调用
	      \end{tasks}
	      \begin{bluebox}{答案与深度解析}
		      \textbf{答案：D. 构造方法在通过 \code{new} 关键字实例化对象时自动调用}
		      \textbf{【核心认知框架】}：考查构造方法语法规则。
		      \textbf{【深度解析】}：
		      \begin{itemize}[leftmargin=*, nosep]
			      \item \textbf{A 错}：构造方法名必须与类名完全相同。
			      \item \textbf{B 错}：构造方法不能有返回值类型（包括 void），否则变为普通方法。
			      \item \textbf{C 错}：支持重载，可以有多个构造方法。
			      \item \textbf{D 对}：\code{new} 指令触发构造方法调用 (\code{invokespecial})。
		      \end{itemize}
	      \end{bluebox}

	\item \textbf{若要在子类的构造方法中调用父类的有参构造方法，必须使用（ \quad ）关键字，并且必须放在子类构造方法的第一行。}
	      \begin{tasks}(4)
		      \task this \task super \task final \task static
	      \end{tasks}
	      \begin{bluebox}{答案与深度解析}
		      \textbf{答案：B. super}
		      \textbf{【核心认知框架】}：考查继承链初始化顺序。
		      \textbf{【深度解析】}：
		      \begin{itemize}[leftmargin=*, nosep]
			      \item \textbf{机制}：\code{super(...)} 调用父类构造器。
			      \item \textbf{限制}：必须在第一行，确保父类先初始化。
			      \item \textbf{默认行为}：若未写，编译器默认插入 \code{super()}。
		      \end{itemize}
	      \end{bluebox}

	\item \textbf{下列关于 \code{final} 关键字的描述，错误的是（ \quad ）。}
	      \begin{tasks}(4)
		      \task 被 \code{final} 修饰的类不能被继承
		      \task 被 \code{final} 修饰的方法不能被重写
		      \task 被 \code{final} 修饰的变量即为常量，一旦赋值后不能修改
		      \task 抽象类可以被 \code{final} 修饰
	      \end{tasks}
	      \begin{bluebox}{答案与深度解析}
		      \textbf{答案：D. 抽象类可以被 \code{final} 修饰}
		      \textbf{【核心认知框架】}：考查 final 与 abstract 的语义冲突。
		      \textbf{【深度解析】}：
		      \begin{itemize}[leftmargin=*, nosep]
			      \item \textbf{冲突}：\code{abstract} 要求被继承，\code{final} 禁止继承，语义矛盾。
			      \item \textbf{编译错误}：\code{public final abstract class A} → 报错。
			      \item \textbf{其他选项}：A、B、C 均为 final 的正确特性。
		      \end{itemize}
	      \end{bluebox}

	\item \textbf{在 Java 中，所有类的顶级父类是（ \quad ）。}
	      \begin{tasks}(4)
		      \task \code{java.lang.Class} \task \code{java.lang.Object} \task \code{java.util.Scanner} \task \code{java.lang.System}
	      \end{tasks}
	      \begin{bluebox}{答案与深度解析}
		      \textbf{答案：B. \code{java.lang.Object}}
		      \textbf{【核心认知框架】}：考查类继承树根节点。
		      \textbf{【深度解析】}：
		      \begin{itemize}[leftmargin=*, nosep]
			      \item \textbf{规范}：JLS §4.3.2 规定所有类隐式继承 Object。
			      \item \textbf{方法}：\code{toString()}, \code{equals()}, \code{hashCode()} 等均源自 Object。
		      \end{itemize}
	      \end{bluebox}

	\item \textbf{接口（Interface）中的成员变量默认的修饰符是（ \quad ）。}
	      \begin{tasks}(4)
		      \task \code{private static final} \task \code{public static final} \task \code{protected static} \task 默认无修饰符
	      \end{tasks}
	      \begin{bluebox}{答案与深度解析}
		      \textbf{答案：B. \code{public static final}}
		      \textbf{【核心认知框架】}：考查接口字段隐式修饰符。
		      \textbf{【深度解析】}：
		      \begin{itemize}[leftmargin=*, nosep]
			      \item \textbf{规范}：JLS §9.3 规定接口字段隐式为 public static final。
			      \item \textbf{必须初始化}：接口字段必须在声明时赋值。
		      \end{itemize}
	      \end{bluebox}

	\item \textbf{当尝试访问数组的非法索引时，Java 运行时会抛出下列哪个异常？（ \quad ）}
	      \begin{tasks}(4)
		      \task \code{NullPointerException} \task \code{ArithmeticException} \task \code{ArrayIndexOutOfBoundsException} \task \code{NumberFormatException}
	      \end{tasks}
	      \begin{bluebox}{答案与深度解析}
		      \textbf{答案：C. \code{ArrayIndexOutOfBoundsException}}
		      \textbf{【核心认知框架】}：考查运行时异常类型。
		      \textbf{【深度解析】}：
		      \begin{itemize}[leftmargin=*, nosep]
			      \item \textbf{机制}：JVM 在数组访问指令 (\code{iaload} 等) 前插入边界检查。
			      \item \textbf{触发}：索引 < 0 或 索引 >= length。
		      \end{itemize}
	      \end{bluebox}

	\item \textbf{下列哪个类主要用于处理频繁修改的字符串序列，且是线程安全的？（ \quad ）}
	      \begin{tasks}(4)
		      \task \code{String} \task \code{StringBuilder} \task \code{StringBuffer} \task \code{Character}
	      \end{tasks}
	      \begin{bluebox}{答案与深度解析}
		      \textbf{答案：C. \code{StringBuffer}}
		      \textbf{【核心认知框架】}：考查字符串类线程安全性区别。
		      \textbf{【深度解析】}：
		      \begin{itemize}[leftmargin=*, nosep]
			      \item \textbf{String}：不可变。
			      \item \textbf{StringBuilder}：可变，非线程安全（快）。
			      \item \textbf{StringBuffer}：可变，线程安全（方法带 synchronized，慢）。
		      \end{itemize}
	      \end{bluebox}

	\item \textbf{Java 集合框架中，如果希望元素能够自动保持排序（如按自然顺序），应选择下列哪个实现类？（ \quad ）}
	      \begin{tasks}(4)
		      \task \code{HashSet} \task \code{ArrayList} \task \code{LinkedList} \task \code{TreeSet}
	      \end{tasks}
	      \begin{bluebox}{答案与深度解析}
		      \textbf{答案：D. \code{TreeSet}}
		      \textbf{【核心认知框架】}：考查集合排序特性。
		      \textbf{【深度解析】}：
		      \begin{itemize}[leftmargin=*, nosep]
			      \item \textbf{TreeSet}：基于 TreeMap (红黑树)，元素有序。
			      \item \textbf{HashSet}：无序。
			      \item \textbf{List}：有序但指插入顺序，非排序。
		      \end{itemize}
	      \end{bluebox}

	\item \textbf{在 Java I/O 流中，以字符为单位进行读写操作的基础抽象类是（ \quad ）。}
	      \begin{tasks}(4)
		      \task \code{InputStream} 和 \code{OutputStream}
		      \task \code{Reader} 和 \code{Writer}
		      \task \code{FileInputStream} 和 \code{FileOutputStream}
		      \task \code{DataInputStream} 和 \code{DataOutputStream}
	      \end{tasks}
	      \begin{bluebox}{答案与深度解析}
		      \textbf{答案：B. \code{Reader} 和 \code{Writer}}
		      \textbf{【核心认知框架】}：考查 IO 流分类（字节 vs 字符）。
		      \textbf{【深度解析】}：
		      \begin{itemize}[leftmargin=*, nosep]
			      \item \textbf{字节流}：InputStream/OutputStream (处理二进制)。
			      \item \textbf{字符流}：Reader/Writer (处理文本，含编码转换)。
		      \end{itemize}
	      \end{bluebox}

	\item \textbf{下列创建和启动多线程的语句中，正确的是（ \quad ）。}
	      \begin{tasks}(4)
		      \task \code{new Thread(...).run();}
		      \task \code{new Thread(...).start();}
		      \task \code{new RunnableImpl().start();}
		      \task \code{Thread.start(...);}
	      \end{tasks}
	      \begin{bluebox}{答案与深度解析}
		      \textbf{答案：B. \code{new Thread(...).start();}}
		      \textbf{【核心认知框架】}：考查线程启动机制。
		      \textbf{【深度解析】}：
		      \begin{itemize}[leftmargin=*, nosep]
			      \item \textbf{start()}：启动新线程，调用 run()。
			      \item \textbf{run()}：普通方法调用，在当前线程执行。
			      \item \textbf{Runnable}：无 start 方法。
		      \end{itemize}
	      \end{bluebox}

	\item \textbf{要使某个类具备被序列化的能力，该类必须实现（ \quad ）接口。}
	      \begin{tasks}(4)
		      \task \code{Cloneable} \task \code{Runnable} \task \code{Serializable} \task \code{Comparable}
	      \end{tasks}
	      \begin{bluebox}{答案与深度解析}
		      \textbf{答案：C. \code{Serializable}}
		      \textbf{【核心认知框架】}：考查序列化标记接口。
		      \textbf{【深度解析】}：
		      \begin{itemize}[leftmargin=*, nosep]
			      \item \textbf{标记接口}：无方法，仅作为 JVM 序列化机制的标志。
			      \item \textbf{机制}：ObjectOutputStream 检查该接口，无则抛 NotSerializableException。
		      \end{itemize}
	      \end{bluebox}

	\item \textbf{GUI 编程中，\code{JFrame} 的默认布局管理器是（ \quad ）。}
	      \begin{tasks}(4)
		      \task \code{FlowLayout} \task \code{GridLayout} \task \code{BorderLayout} \task \code{CardLayout}
	      \end{tasks}
	      \begin{bluebox}{答案与深度解析}
		      \textbf{答案：C. \code{BorderLayout}}
		      \textbf{【核心认知框架】}：考查 Swing 容器默认布局。
		      \textbf{【深度解析】}：
		      \begin{itemize}[leftmargin=*, nosep]
			      \item \textbf{JFrame}：内容面板默认 BorderLayout (北南西中东)。
			      \item \textbf{JPanel}：默认 FlowLayout。
		      \end{itemize}
	      \end{bluebox}

	\item \textbf{JDBC 操作数据库时，用于执行预编译 SQL 语句的对象是（ \quad ）。}
	      \begin{tasks}(4)
		      \task \code{Statement} \task \code{PreparedStatement} \task \code{CallableStatement} \task \code{ResultSet}
	      \end{tasks}
	      \begin{bluebox}{答案与深度解析}
		      \textbf{答案：B. \code{PreparedStatement}}
		      \textbf{【核心认知框架】}：考查 JDBC 接口区别。
		      \textbf{【深度解析】}：
		      \begin{itemize}[leftmargin=*, nosep]
			      \item \textbf{PreparedStatement}：预编译，防 SQL 注入，效率高。
			      \item \textbf{Statement}：普通执行。
			      \item \textbf{CallableStatement}：存储过程。
		      \end{itemize}
	      \end{bluebox}
\end{enumerate}

\section*{二、判断题深度解析（共 10 小题）}
\textit{（正确的选 A，错误的选 B，深度分析命题真伪）}

\begin{enumerate}[label=\textbf{\arabic*.}, leftmargin=*, itemsep=1.5em]
	\item \textbf{Java 中的数组一旦创建，其长度就固定不可改变。}
	      \begin{bluebox}{答案与深度解析}
		      \textbf{答案：A (正确)}
		      \textbf{【核心认知框架】}：考查数组对象的不可变性。
		      \textbf{【深度解析】}：
		      \begin{itemize}[leftmargin=*, nosep]
			      \item \textbf{内存模型}：数组对象在堆中分配，length 字段为 final。
			      \item \textbf{字节码}：\code{arraylength} 指令读取长度，无写入指令。
			      \item \textbf{对比}：ArrayList 是可变长度的封装，底层数组仍固定（扩容是建新数组）。
		      \end{itemize}
	      \end{bluebox}

	\item \textbf{\code{==} 运算符用于比较两个对象的内存地址是否相同，而 \code{equals()} 方法在未被重写时比较的是对象的属性值。}
	      \begin{bluebox}{答案与深度解析}
		      \textbf{答案：B (错误)}
		      \textbf{【核心认知框架】}：考查 equals 默认行为。
		      \textbf{【深度解析】}：
		      \begin{itemize}[leftmargin=*, nosep]
			      \item \textbf{真相}：Object 类的 \code{equals()} 默认实现就是 \code{this == obj}，比较地址。
			      \item \textbf{重写}：只有重写后（如 String）才比较内容。
			      \item \textbf{命题陷阱}：混淆“默认行为”与“重写后行为”。
		      \end{itemize}
	      \end{bluebox}

	\item \textbf{在 Java 中，一个子类只能继承一个父类，但一个类可以实现多个接口。}
	      \begin{bluebox}{答案与深度解析}
		      \textbf{答案：A (正确)}
		      \textbf{【核心认知框架】}：考查单继承多实现机制。
		      \textbf{【深度解析】}：
		      \begin{itemize}[leftmargin=*, nosep]
			      \item \textbf{设计原因}：避免菱形继承问题（C++ 痛点）。
			      \item \textbf{字节码}：Class 文件结构中 only one super\_class，multiple interfaces 数组。
		      \end{itemize}
	      \end{bluebox}

	\item \textbf{父类的构造方法可以被子类继承和直接使用。}
	      \begin{bluebox}{答案与深度解析}
		      \textbf{答案：B (错误)}
		      \textbf{【核心认知框架】}：考查构造方法继承性。
		      \textbf{【深度解析】}：
		      \begin{itemize}[leftmargin=*, nosep]
			      \item \textbf{规则}：构造方法不可继承。
			      \item \textbf{调用}：子类通过 \code{super()} 隐式或显式调用。
		      \end{itemize}
	      \end{bluebox}

	\item \textbf{带有 \code{try-catch-finally} 的异常处理结构中，无论是否发生异常，\code{finally} 块中的代码通常都会被执行。}
	      \begin{bluebox}{答案与深度解析}
		      \textbf{答案：A (正确)}
		      \textbf{【核心认知框架】}：考查 finally 执行语义。
		      \textbf{【深度解析】}：
		      \begin{itemize}[leftmargin=*, nosep]
			      \item \textbf{例外}：\code{System.exit(0)} 或 JVM 崩溃时不执行。
			      \item \textbf{字节码}：\code{jsr} (旧) 或异常表属性保证 finally 跳转。
		      \end{itemize}
	      \end{bluebox}

	\item \textbf{\code{String} 类型的对象是不可变的（Immutable），一旦创建就不能修改其内容。}
	      \begin{bluebox}{答案与深度解析}
		      \textbf{答案：A (正确)}
		      \textbf{【核心认知框架】}：考查 String 不可变性。
		      \textbf{【深度解析】}：
		      \begin{itemize}[leftmargin=*, nosep]
			      \item \textbf{底层}：char[] 或 byte[] 被 final 修饰且无修改方法。
			      \item \textbf{优势}：线程安全、哈希码缓存、字符串池安全。
		      \end{itemize}
	      \end{bluebox}

	\item \textbf{在多重 \code{catch} 块中，必须将捕获 \code{Exception} 类的块放在捕获 \code{IOException} 类的块之前。}
	      \begin{bluebox}{答案与深度解析}
		      \textbf{答案：B (错误)}
		      \textbf{【核心认知框架】}：考查 catch 块顺序规则。
		      \textbf{【深度解析】}：
		      \begin{itemize}[leftmargin=*, nosep]
			      \item \textbf{规则}：子类异常必须放在父类异常之前。
			      \item \textbf{原因}：否则子类 catch 块永远无法到达（Unreachable Code）。
			      \item \textbf{正确}：\code{IOException} 在前，\code{Exception} 在后。
		      \end{itemize}
	      \end{bluebox}

	\item \textbf{\code{ArrayList} 内部是基于动态数组实现的，因此它的随机访问速度比 \code{LinkedList} 更快。}
	      \begin{bluebox}{答案与深度解析}
		      \textbf{答案：A (正确)}
		      \textbf{【核心认知框架】}：考查集合底层数据结构性能。
		      \textbf{【深度解析】}：
		      \begin{itemize}[leftmargin=*, nosep]
			      \item \textbf{ArrayList}：O(1) 随机访问 (get)。
			      \item \textbf{LinkedList}：O(n) 随机访问 (需遍历)。
		      \end{itemize}
	      \end{bluebox}

	\item \textbf{包含抽象方法的类必须声明为抽象类，但抽象类也可以不包含任何抽象方法。}
	      \begin{bluebox}{答案与深度解析}
		      \textbf{答案：A (正确)}
		      \textbf{【核心认知框架】}：考查抽象类定义规则。
		      \textbf{【深度解析】}：
		      \begin{itemize}[leftmargin=*, nosep]
			      \item \textbf{强制}：有抽象方法 → 必须是抽象类。
			      \item \textbf{允许}：抽象类 → 可以无抽象方法（用于限制实例化）。
		      \end{itemize}
	      \end{bluebox}

	\item \textbf{方法内部声明的局部变量如果没有进行显式初始化，系统会自动为其赋予默认值。}
	      \begin{bluebox}{答案与深度解析}
		      \textbf{答案：B (错误)}
		      \textbf{【核心认知框架】}：考查局部变量初始化规则。
		      \textbf{【深度解析】}：
		      \begin{itemize}[leftmargin=*, nosep]
			      \item \textbf{规则}：局部变量无默认值，使用前必须初始化。
			      \item \textbf{对比}：成员变量有默认值 (0, null, false)。
			      \item \textbf{编译错误}：\code{variable might not have been initialized}。
		      \end{itemize}
	      \end{bluebox}
\end{enumerate}

\section*{三、填空题深度解析（共 5 小题）}

\begin{enumerate}[label=\textbf{\arabic*.}, leftmargin=*, itemsep=1.5em]
	\item \textbf{Java 的跨平台特性主要依赖于 Java 虚拟机（JVM）。Java 源代码被编译后生成的中间文件称为\underline{\hspace{4cm}}，其扩展名是\underline{\hspace{3cm}}。}
	      \begin{bluebox}{答案与深度解析}
		      \textbf{答案：字节码（或字节码文件）；.class}
		      \textbf{【核心认知框架】}：考查 Java 编译产物。
		      \textbf{【深度解析】}：
		      \begin{itemize}[leftmargin=*, nosep]
			      \item \textbf{机制}：javac 编译源文件为字节码，JVM 解释执行。
			      \item \textbf{跨平台}：字节码统一，JVM 适配各 OS。
		      \end{itemize}
	      \end{bluebox}

	\item \textbf{面向对象的三大核心特征分别是：封装、\underline{\hspace{4cm}} 和 \underline{\hspace{4cm}}。}
	      \begin{bluebox}{答案与深度解析}
		      \textbf{答案：继承；多态}
		      \textbf{【核心认知框架】}：考查 OOP 理论基础。
		      \textbf{【深度解析】}：
		      \begin{itemize}[leftmargin=*, nosep]
			      \item \textbf{封装}：隐藏细节，暴露接口。
			      \item \textbf{继承}：代码复用，is-a 关系。
			      \item \textbf{多态}：同一接口不同实现，动态绑定。
		      \end{itemize}
	      \end{bluebox}

	\item \textbf{在 Java 的异常处理体系中，所有的异常类都继承自 \code{Throwable} 类，它主要有两个直接子类：\underline{\hspace{4cm}}（用于系统严重错误）和 \underline{\hspace{4cm}}（用于程序可以处理的异常）。}
	      \begin{bluebox}{答案与深度解析}
		      \textbf{答案：Error；Exception}
		      \textbf{【核心认知框架】}：考查异常类继承树。
		      \textbf{【深度解析】}：
		      \begin{itemize}[leftmargin=*, nosep]
			      \item \textbf{Error}：JVM 层面错误（如 OOM），不可恢复。
			      \item \textbf{Exception}：程序逻辑错误，可捕获处理。
		      \end{itemize}
	      \end{bluebox}

	\item \textbf{Java 中线程的生命周期包含了新建（New）、就绪（Runnable）、\underline{\hspace{3cm}}（Running）、\underline{\hspace{3cm}}（Blocked）以及死亡（Dead）五种基本状态。}
	      \begin{bluebox}{答案与深度解析}
		      \textbf{答案：运行（状态）；阻塞（状态）}
		      \textbf{【核心认知框架】}：考查线程状态模型。
		      \textbf{【深度解析】}：
		      \begin{itemize}[leftmargin=*, nosep]
			      \item \textbf{JVM 规范}：JVMS §2.7 定义线程状态。
			      \item \textbf{API 映射}：Thread.State 枚举 (NEW, RUNNABLE, BLOCKED, WAITING, TIMED\_WAITING, TERMINATED)。
		      \end{itemize}
	      \end{bluebox}

	\item \textbf{在 Java 集合框架中，\underline{\hspace{4cm}} 接口的实现类允许存在重复的元素并且是有序的；而 \underline{\hspace{4cm}} 接口的实现类不允许存在重复的元素。}
	      \begin{bluebox}{答案与深度解析}
		      \textbf{答案：List；Set}
		      \textbf{【核心认知框架】}：考查 Collection 子接口特性。
		      \textbf{【深度解析】}：
		      \begin{itemize}[leftmargin=*, nosep]
			      \item \textbf{List}：有序可重复 (ArrayList, LinkedList)。
			      \item \textbf{Set}：无序不可重复 (HashSet, TreeSet)。
		      \end{itemize}
	      \end{bluebox}
\end{enumerate}

\section*{四、程序运行结果题深度解析（共 5 小题）}

\begin{enumerate}[label=\textbf{\arabic*.}, leftmargin=*, itemsep=1.5em]
	\item \textbf{写出下列程序的输出结果：}
	      \begin{lstlisting}[language=Java]
public class Test1 {
public static void main(String[] args) {
int sum = 0;
for (int i = 1; i <= 5; i++) {
if (i == 3) continue;
sum += i;
}
System.out.println("sum = " + sum);
}
}
\end{lstlisting}
	      \begin{bluebox}{答案与深度解析}
		      \textbf{答案：\code{sum = 12}}
		      \textbf{【核心认知框架】}：考查循环控制流 continue 语义。
		      \textbf{【执行流程深度拆解】}：
		      \begin{itemize}[leftmargin=*, nosep]
			      \item \textbf{逻辑}：i=1,2,4,5 累加；i=3 跳过。
			      \item \textbf{计算}：1+2+4+5 = 12。
			      \item \textbf{字节码}：\code{continue} 对应跳转指令跳过累加步骤。
		      \end{itemize}
	      \end{bluebox}

	\item \textbf{写出下列程序的输出结果：}
	      \begin{lstlisting}[language=Java]
public class Test2 {
public static void change(int[] arr) {
arr[0] = 99;
}
public static void main(String[] args) {
int[] nums = {10, 20, 30};
change(nums);
System.out.println(nums[0] + ", " + nums[1]);
}
}
\end{lstlisting}
	      \begin{bluebox}{答案与深度解析}
		      \textbf{答案：\code{99, 20}}
		      \textbf{【核心认知框架】}：考查数组引用传递机制。
		      \textbf{【执行流程深度拆解】}：
		      \begin{itemize}[leftmargin=*, nosep]
			      \item \textbf{传递}：传递数组引用副本，指向同一堆对象。
			      \item \textbf{修改}：方法内修改堆内容，外部可见。
			      \item \textbf{结果}：nums[0] 变为 99，nums[1] 不变。
		      \end{itemize}
	      \end{bluebox}

	\item \textbf{写出下列程序的输出结果：}
	      \begin{lstlisting}[language=Java]
class Father {
int x = 10;
public void show() { System.out.println("Father Show"); }
}
class Son extends Father {
int x = 20;
public void show() { System.out.println("Son Show"); }
}
public class Test3 {
public static void main(String[] args) {
Father f = new Son();
System.out.println(f.x);
f.show();
}
}
\end{lstlisting}
	      \begin{bluebox}{答案与深度解析}
		      \textbf{答案：\code{10}} \newline \code{Son Show}
		      \textbf{【核心认知框架】}：考查多态与方法/字段绑定区别。
		      \textbf{【执行流程深度拆解】}：
		      \begin{itemize}[leftmargin=*, nosep]
			      \item \textbf{字段}：非多态，看左边（父类），输出 10。
			      \item \textbf{方法}：多态，看右边（子类），输出 Son Show。
			      \item \textbf{字节码}：字段访问 \code{getfield} 静态绑定，方法调用 \code{invokevirtual} 动态绑定。
		      \end{itemize}
	      \end{bluebox}

	\item \textbf{写出下列程序的输出结果：}
	      \begin{lstlisting}[language=Java]
public class Test4 {
public static void main(String[] args) {
String s1 = new String("Java");
String s2 = new String("Java");
System.out.println(s1 == s2);
System.out.println(s1.equals(s2));
}
}
\end{lstlisting}
	      \begin{bluebox}{答案与深度解析}
		      \textbf{答案：\code{false}} \newline \code{true}
		      \textbf{【核心认知框架】}：考查字符串池与 equals 重写。
		      \textbf{【执行流程深度拆解】}：
		      \begin{itemize}[leftmargin=*, nosep]
			      \item \textbf{new String}：每次在堆上新建对象，地址不同，\code{==} 为 false。
			      \item \textbf{equals}：String 重写为比较内容，内容相同为 true。
		      \end{itemize}
	      \end{bluebox}

	\item \textbf{写出下列程序的输出结果：}
	      \begin{lstlisting}[language=Java]
public class Test5 {
public static int testException() {
int res = 1;
try { res = 10 / 0; return res; } 
catch (ArithmeticException e) { res = 2; return res; } 
finally { res = 3; System.out.print("Finally "); }
}
public static void main(String[] args) {
System.out.println(testException());
}
}
\end{lstlisting}
	      \begin{bluebox}{答案与深度解析}
		      \textbf{答案：\code{Finally 2}}
		      \textbf{【核心认知框架】}：考查 finally 与 return 执行顺序。
		      \textbf{【执行流程深度拆解】}：
		      \begin{itemize}[leftmargin=*, nosep]
			      \item \textbf{流程}：try 异常 → catch 捕获 → res=2 → 准备返回 → 执行 finally → 输出 Finally → 返回 catch 中暂存的 2。
			      \item \textbf{陷阱}：finally 中修改 res 不影响返回值（基本类型），除非 finally 中有 return。
		      \end{itemize}
	      \end{bluebox}
\end{enumerate}

\section*{五、编程题深度解析与设计模式分析（共 2 小题）}

\begin{enumerate}[label=\textbf{\arabic*.}, leftmargin=*, itemsep=1.5em]
	\item \textbf{面向对象基础设计（银行账户类）}
	      \begin{bluebox}{答案与深度解析}
		      \textbf{核心设计思路：封装性 + 业务逻辑校验}

		      \textbf{【核心认知框架】}
		      \begin{itemize}[leftmargin=*, nosep]
			      \item \textbf{题目本质}：考查类的基本结构、封装性及方法逻辑。
			      \item \textbf{关键区分点}：私有属性、构造器初始化、业务规则（余额检查）。
			      \item \textbf{命题人意图}：测试 POJO 设计规范及基本逻辑控制。
		      \end{itemize}

		      \textbf{【关键技术点深度拆解】}
		      \begin{itemize}[leftmargin=*, nosep]
			      \item \textbf{封装性}：
			            \begin{itemize}
				            \item 属性 \code{private}，防止外部直接修改余额。
				            \item 通过 \code{deposit}/\code{withdraw} 方法控制状态变更。
			            \end{itemize}
			      \item \textbf{构造方法}：
			            \begin{itemize}
				            \item 使用 \code{this} 区分参数与成员变量。
				            \item 确保对象创建时状态合法。
			            \end{itemize}
			      \item \textbf{业务逻辑}：
			            \begin{itemize}
				            \item 取款前检查 \code{balance >= amount}。
				            \item 体现面向对象的行为建模。
			            \end{itemize}
		      \end{itemize}

		      \textbf{【参考代码核心片段】}
		      \begin{lstlisting}[language=Java]
public void withdraw(double amount) {
    if (balance >= amount) {
        balance -= amount;
        System.out.println("取款成功");
    } else {
        System.out.println("余额不足");
    }
}
\end{lstlisting}
	      \end{bluebox}

	\item \textbf{接口与多态的综合应用（图形面积计算）}
	      \begin{bluebox}{答案与深度解析}
		      \textbf{核心设计思路：接口隔离 + 多态数组}

		      \textbf{【核心认知框架】}
		      \begin{itemize}[leftmargin=*, nosep]
			      \item \textbf{题目本质}：考查接口定义、实现及多态调用。
			      \item \textbf{关键区分点}：接口引用指向实现类对象，统一调用 \code{getArea()}。
			      \item \textbf{命题人意图}：测试多态性在解耦中的实际应用。
		      \end{itemize}

		      \textbf{【关键技术点深度拆解】}
		      \begin{itemize}[leftmargin=*, nosep]
			      \item \textbf{接口定义}：
			            \begin{itemize}
				            \item \code{Shape} 定义行为契约 \code{getArea()}。
				            \item 符合依赖倒置原则（DIP）。
			            \end{itemize}
			      \item \textbf{多态数组}：
			            \begin{itemize}
				            \item \code{Shape[] shapes} 可存放任何实现类。
				            \item 循环调用 \code{shapes[i].getArea()} 动态绑定。
			            \end{itemize}
			      \item \textbf{扩展性}：
			            \begin{itemize}
				            \item 新增图形只需实现接口，无需修改测试代码。
			            \end{itemize}
		      \end{itemize}

		      \textbf{【参考代码核心片段】}
		      \begin{lstlisting}[language=Java]
Shape[] shapes = new Shape[2];
shapes[0] = new Rectangle(4, 5);
shapes[1] = new Circle(3);
// 多态调用
for (Shape s : shapes) {
    totalArea += s.getArea();
}
\end{lstlisting}
	      \end{bluebox}
\end{enumerate}

\vfill
\begin{center}
	\zihao{4}\bfseries —— 讲义结束，祝学习顺利 ——
\end{center}
\end{document}
